\section{全范围拼接与六次幂结论}
\label{sec:assembly}

本节将各段的实际素数定理拼成关于所有正整数的结论。
需要区分两个任务：前面各节已经证明具体范围中的素数存在性；本节再证明
这些范围覆盖全部正整数，并核对实数幂、整数幂及所有严格端点。
这一步没有再引入零点估计或数值假设。

为与形式化接口一致，记
\begin{align}
 \mathcal S_6(x)&:\quad
   \exists p\in\N,\quad p\text{ 为素数},\quad x<p<x+6x^{5/6},
   \label{assembly:S6}\\
 \mathcal S_4(x)&:\quad
   \exists p\in\N,\quad p\text{ 为素数},\quad x<p<x+4x^{3/4},
   \label{assembly:S4}\\
 \mathcal N_6(n)&:\quad
   \exists p\in\N,\quad p\text{ 为素数},\quad n^6<p<n^6+6n^5.
   \label{assembly:N6}
\end{align}
前两个定义中的比较在 $\R$ 中进行，自然数 $p$ 通过标准嵌入视为实数；
第三个定义中的幂与不等式都是自然数运算。
对应定义分别是 \leanref{Definitions}{ShortInterval6}、
\leanref{Definitions}{ShortInterval4} 和
\leanref{Definitions}{NatSixthShort}。
下标 $4$ 和 $6$ 在这里表示选用的短区间长度；
它们本身不是某个全范围幂次猜想已经成立的声明。

\subsection{已经获得的各段及其精确依赖}

所有实变量段都以 $x>0$、$L=\log x$ 为前提。
下表把用于最终拼接的五个结论放在一起。
其中的每个区间端点均包含在相应定理中。

\begin{center}
\small
\begin{tabular}{c|c|c|l}
范围 & 变量及条件 & 结论 & 外部输入\\\hline
整数下段 & $n\geq1$, $\log(n^6)\leq167$ & $\mathcal N_6(n)$ & H7\\
下中段 & $167\leq L\leq715/4$ & $\mathcal S_6(x)$ & H1, H2, H3, H4, H5\\
上中段 & $715/4\leq L\leq400$ & $\mathcal S_6(x)$ & H1, H2, H3, H4\\
有限远段 & $400\leq L\leq700$ & $\mathcal S_4(x)$ & H1, H2, H3, H4\\
无限尾部 & $700\leq L$ & $\mathcal S_4(x)$ & H1, H2, H3, H6
\end{tabular}
\end{center}

整数下段是第~\ref{sec:lower} 节的定理~\ref{lower:log-theorem}。
下中段与上中段分别由
\leanref{M3LowerResult}{lower_middle_short_interval_from_inputs} 和
\leanref{M3MiddleUpper}{middle_upper_short_interval_from_inputs} 给出。
有限远段为
\leanref{M3FarUpper}{far_short_interval_from_inputs}，
无限尾部为第~\ref{sec:tail} 节的定理~\ref{tail:prime-theorem}。
中段分点 $715/4=178.75$ 只决定采用哪一个已证明的解析估计，
不会在最终定理中留下这个分点处的额外假设。

以上依赖表也说明，H5 的使用仅在下中段，H6 仅在无限尾部，H7 仅在整数下段。
最终定理统一接受七个输入，是因为它同时覆盖这些范围；
这不表示每一个局部定理都使用了全部七个输入。

\subsection{四次长度与六次长度的比较}

\begin{lemma}[短区间长度的包含关系]\label{assembly:length-comparison}
若 $x\geq1$，则
\begin{equation}\label{assembly:length-inequality}
 4x^{3/4}\leq6x^{5/6}.
\end{equation}
因此 $\mathcal S_4(x)$ 蕴含 $\mathcal S_6(x)$。
\end{lemma}

\begin{proof}
因为 $3/4\leq5/6$ 且 $x\geq1$，实数幂关于指数的单调性给出
$x^{3/4}\leq x^{5/6}$。后者非负，故
\[
 4x^{3/4}\leq4x^{5/6}\leq6x^{5/6}.
\]
若 $p$ 是 $\mathcal S_4(x)$ 的素数见证，则
\[
 x<p<x+4x^{3/4}\leq x+6x^{5/6}.
\]
严格不等式与非严格不等式的传递仍给出严格的最终右界。
\end{proof}

相应实现为 \leanref{Reduction}{four_length_le_six_length} 和
\leanref{Reduction}{short_six_of_short_four}。
这里必须说明 $x\geq1$；当 $0<x<1$ 时，单凭指数大小不能作同方向比较。
在本证明的有限远段和无限尾部，$L\geq400$，所以这个条件自动满足；
在最终应用 $x=n^6$ 中，它也由 $n\geq1$ 直接得到。

\begin{proposition}[统一的解析上段]\label{assembly:analytic-range}
在 H1 至 H6 下，对每个 $x>0$，若 $\log x\geq167$，则
$\mathcal S_6(x)$ 成立。
\end{proposition}

\begin{proof}
记 $L=\log x$。若 $L\leq400$，则在
$L\leq715/4$ 与 $L>715/4$ 两种情况下，分别使用下中段与上中段的
$\mathcal S_6(x)$ 结论。
若 $L>400$，则 $x>1$；当 $L\leq700$ 时使用有限远段，
当 $L>700$ 时使用无限尾部。后两种情况先得到 $\mathcal S_4(x)$，
再由引理~\ref{assembly:length-comparison} 得到 $\mathcal S_6(x)$。
每种情形均由某个闭区间定理覆盖，故 $167$、$715/4$、$400$、$700$
这四个接口没有遗漏。
\end{proof}

中段的具体合并见
\leanref{FinalResult}{middle_short_interval_from_inputs}。
把三个已构造的解析段交给
\leanref{Assembly}{real_short_above_167_of_segment_inputs}
便得到命题~\ref{assembly:analytic-range} 的形式化版本。

\subsection{从实变量区间回到整数六次幂}

\begin{lemma}[实数幂与自然数幂的精确转换]\label{assembly:power-conversion}
对每个自然数 $n$，有
\begin{equation}\label{assembly:rpow-identity}
 \bigl((n:\R)^6\bigr)^{5/6}=(n:\R)^5.
\end{equation}
若 $\mathcal S_6((n:\R)^6)$ 成立，则 $\mathcal N_6(n)$ 成立。
\end{lemma}

\begin{proof}
自然数在实数中的像非负，故非负实底数的幂法则适用。
由 $6(5/6)=5$，得到式~\eqref{assembly:rpow-identity}。
将这个等式代入 $\mathcal S_6$ 的右端点后，见证素数满足
\[
 (n:\R)^6<(p:\R)<(n:\R)^6+6(n:\R)^5.
\]
自然数到实数的嵌入保持加法、乘法、整数幂和严格序关系，
因此两条实数不等式恰给出式~\eqref{assembly:N6} 的两条自然数不等式。
素数 $p$ 本身不发生改变。
\end{proof}

实现为 \leanref{Reduction}{sixth_rpow_five_sixths} 与
\leanref{Reduction}{nat_short_of_real_short}。
这一转换没有使用四舍五入，也不需要对实数端点先取整。

\begin{lemma}[相邻六次幂之间的严格余量]\label{assembly:binomial-gap}
对每个自然数 $n$，
\begin{equation}\label{assembly:binomial}
 n^6+6n^5<(n+1)^6.
\end{equation}
\end{lemma}

\begin{proof}
二项式展开给出
\[
 (n+1)^6-(n^6+6n^5)
 =15n^4+20n^3+15n^2+6n+1>0.
\]
各幂非负，而常数项为 $1$，所以严格正性甚至对 $n=0$ 也成立。
\end{proof}

对应定理是 \leanref{Reduction}{sixth_binomial_gap}。
不过最终的素数存在性定理仍只声明 $n\geq1$：$n=0$ 时区间 $(0,1)$
没有素数，不能因为这个单独的二项式比较对 $n=0$ 成立便扩大主定理范围。

\subsection{全范围条件定理}

\begin{theorem}[七个明确输入下的六次幂短区间定理]
\label{assembly:main-short}
假设 H1 至 H7 成立。对每个自然数 $n\geq1$，存在自然数素数 $p$，使
\begin{equation}\label{assembly:strong-chain}
 n^6<p<n^6+6n^5<(n+1)^6.
\end{equation}
\end{theorem}

\begin{proof}
固定 $n\geq1$，令 $x=(n:\R)^6>0$，$L=\log x$。
若 $L\leq167$，由定理~\ref{lower:log-theorem} 直接得到
$\mathcal N_6(n)$。
若 $L>167$，则命题~\ref{assembly:analytic-range} 给出
$\mathcal S_6(x)$，再由引理~\ref{assembly:power-conversion} 得
$\mathcal N_6(n)$。
两种情形已经穷尽所有 $n\geq1$，且各自均给出实际的自然数素数见证。
最后附上引理~\ref{assembly:binomial-gap} 的严格比较，得到完整链式不等式。
\end{proof}

\begin{corollary}[六次幂之间的素数]\label{assembly:sixth-powers}
在 H1 至 H7 下，对每个自然数 $n\geq1$，开区间
$(n^6,(n+1)^6)$ 中至少有一个素数。
\end{corollary}

这两个结论分别由
\leanref{FinalResult}{sixth_power_short_from_named_inputs} 及
\leanref{FinalResult}{sixth_powers_from_named_inputs} 实现。
后一个定理的目标
\leanref{Definitions}{SixthPowerConclusion}
直接量化所有 $n\in\N$ 并要求 $1\leq n$；它不是“所有已检查整数”的缩写。
前一个定理保留了更强的内端点 $n^6+6n^5$，所以同时记录了
式~\eqref{assembly:strong-chain} 中的额外余量。

\subsection{中间接口已经构造，而非留下局部结论作为前提}

证明库保留两个便于复用的中间接口。
\leanref{Assembly}{AnalyticSegmentInputs} 的三个字段分别陈述
中段、有限远段和尾部的素数结论；
\leanref{AssemblyWithTail}{FiniteAnalyticSegmentInputs}
只包含前两个有限段。
如果仅证明“给定这些字段，就能推出六次幂结论”，那仍然只是一个拼接定理，
因为困难的局部素数结论仍作为未证明字段存在。

当前最终定理补齐了这一步。
\leanref{FinalResult}{finite_analytic_segments_from_inputs}
把已证明的中段与有限远段定理填入两个字段；
\leanref{TailFromInputs}{tail_short_interval_from_inputs}
再提供尾部字段。
\leanref{FinalResult}{sixth_power_short_from_named_inputs}
和 \leanref{FinalResult}{sixth_powers_from_named_inputs}
的参数表因此只留下 H1 至 H7，并不再要求调用者提供
一个局部素数结论、一个单元最大值不等式或一个未证明的解析段接口。

这一点还可以在历史单元的对应关系中直接检查。
\leanref{FinalCellCoverage}{all_94_cells_actual_prime_intervals}
证明每个保存的中段细分单元内都有 $\mathcal S_6(x)$；
\leanref{FinalCellCoverage}{all_60_far_cells_actual_prime_intervals}
证明每个保存的远段单元内都有 $\mathcal S_4(x)$。
它们先在内核中检查每个单元落入相应的全段，再调用实际的素数定理。
因此保存的有理数最大值并不是这两个结论的额外解析前提。
同时，当前远段的形式化使用梯形权函数。
它覆盖历史六十个单元的目标区间，但不等于把历史多项式权函数手稿逐行翻译为 Lean。
这两种研究记录应分别保留，避免将“目标范围已形式化”误写成
“每一种历史证明实现均已形式化”。

\begin{remark}[主定理的逻辑范围]\label{assembly:scope}
本文完成的是七个明确数学输入之后的全范围推导。
H1 是本路线承认的 QRH 假设；H2 至 H7 是精确定义的文献命题接口，
其来源状态及独立审查程度须分别陈述。
Lean 证明这些命题一旦成立，式~\eqref{assembly:strong-chain}
便对每个 $n\geq1$ 成立；它不凭编译成功而证明七个输入本身。
特别地，这不是只以 QRH 一项为参数的库内定理，也不是无条件六次幂定理。
至于 $\mathcal S_4(x)$，本文已经在 $\log x\geq400$ 的实变量范围证明它，
但没有据此宣称所有正整数的四次幂猜想已经解决。
\end{remark}
